\documentclass[10pt,a4paper]{amsart}
\usepackage[margin=19mm]{geometry}
\usepackage{amsmath,amssymb,mathtools}
\usepackage[scr=rsfs,cal=euler]{mathalfa}
\usepackage{xfrac}
\usepackage[unicode,hidelinks,bookmarks=false]{hyperref}
\newcommand{\ie}{i.e.\ }
\newcommand{\cf}{cf.\ }
\newcommand{\resp}{respectively\ }
\newcommand{\Subsection}{\subsection}
\providecommand{\nobreakdash}{\nobreak\hbox{-}}
\setlength{\emergencystretch}{2em}
\multlinegap0pt
\usepackage{amssymb}
\usepackage{tikz}
\usepackage{bm}
\usepackage{enumerate}
\usepackage[shortlabels]{enumitem}
\usepackage{csquotes}
\usepackage{float}
\usepackage{xcolor}
\usepackage{mathtools}
\newtheorem{lemma}{Lemma}
\title{The sharp asymptotic density of zero-sum-free spherical sets}
\author{Hong-Jun Ge, Zixiang Xu}
\date{Source: arXiv:2607.05099v2 (CC BY 4.0)}
\begin{document}
\maketitle

\noindent\emph{Source and licence.} The sharp asymptotic density of zero-sum-free spherical sets. Version arXiv:2607.05099v2 (CC BY 4.0), distributed by arXiv under the Creative Commons Attribution 4.0 International licence, http://creativecommons.org/licenses/by/4.0/, as recorded in the arXiv licence metadata for this exact version and shown on the abstract page https://arxiv.org/abs/2607.05099v2. Full licence notice: \texttt{licenses/arxiv-2607.05099-CC-BY-4.0.txt}.\par\smallskip

\noindent\textit{Editorial scope.} Counted unit: the article's lemma lem:app-H4-A4-incidences (source0.tex lines 901--920). The article's complete original proof of it is delivered with it (source0.tex lines 922--1027, one \texttt{proof} environment of 106 lines). No mathematics is written, repaired or re-derived: the statement and the proof are the article's own lines, in the article's own notation, and no retained line is substituted or re-flowed.  Retained uncounted as the source-local context the proof needs: lines 799--811.  Nothing was omitted from any retained span, so the package carries no omission record.  No retained line contains a citation, so the wrapper carries no bibliography and no bibliography entry.  No retained line contains a figure environment, an image-inclusion command or drawing code, so no image is delivered and the package declares no asset dependency.  Editorial formatting only, in addition: an AMS \texttt{amsart} preamble that restates the article's own declarations and macros used by the retained lines, namely the environments \texttt{lemma} on separate counters, together with the standard packages those lines need.  The retained environments print in this document as Lemma 1 in order of appearance, whereas the article prints the same items with the numbering of its own full document.  The complete native source of the article as distributed in the arXiv e-print for this exact version is bundled unchanged as \texttt{sources/204514/source0.tex}.
\begin{lemma}\label{lem:app-H4-line-valencies}
For every line \(L\in\mathcal L\), the other \(59\) lines split as follows:
\begin{equation}\label{eq:app-H4-line-valencies}
\begin{array}{c|cccc}
|\langle L,M\rangle|
& \frac{1}{2} & 0 & \frac{\varphi^{-1}}{2} & \frac{\varphi}{2} \\
\hline
\#\{M:|\langle L,M\rangle|\text{ has this value}\}
& 20 & 15 & 12 & 12 .
\end{array}
\end{equation}
In particular, \(\Gamma\) is \(20\)-regular.
\end{lemma}

\begin{lemma}\label{lem:app-H4-A4-incidences}
The family \(\mathcal C\) of \(A_{4}\)-subsystems satisfies the following.
\begin{enumerate}
    \item[\textup{(i)}] \(|\mathcal C|=60\).
    \item[\textup{(ii)}] Every line \(L\in\mathcal L\) belongs to exactly
    \(10\) members of \(\mathcal C\).
    \item[\textup{(iii)}] For two lines \(L,M\in\mathcal L\), the number
    of members of \(\mathcal C\) containing both \(L\) and \(M\) depends
    only on \(|\langle L,M\rangle|\), and is given by
    \begin{equation}\label{eq:app-H4-A4-pair-table}
    \begin{array}{c|ccccc}
    |\langle L,M\rangle|
    & 1 & \frac{1}{2} & 0 & \frac{\varphi^{-1}}{2} & \frac{\varphi}{2} \\
    \hline
    \#\{C\in\mathcal C:L,M\in C\}
    & 10 & 3 & 2 & 0 & 0 .
    \end{array}
    \end{equation}
\end{enumerate}
\end{lemma}

\begin{proof}[Proof of Lemma~\ref{lem:app-H4-A4-incidences}]
We first display the construction for
\(L_{0}=\langle(1,0,0,0)\rangle\).  Let \(\mathcal N_{0}(L_{0})\) be the
set of lines orthogonal to \(L_{0}\).  By
Lemma~\ref{lem:app-H4-line-valencies},
\(|\mathcal N_{0}(L_{0})|=15\).  Write these lines as
\[
    E_{1}=\langle(0,1,0,0)\rangle,\ 
    E_{2}=\langle(0,0,1,0)\rangle,\ 
    E_{3}=\langle(0,0,0,1)\rangle,
\]
and, for \(\epsilon,\eta\in\{\pm1\}\),
\[
    X_{\epsilon,\eta}=\langle(0,a,\epsilon b,\eta c)\rangle,\ 
    Y_{\epsilon,\eta}=\langle(0,c,\epsilon a,\eta b)\rangle,\ 
    Z_{\epsilon,\eta}=\langle(0,b,\epsilon c,\eta a)\rangle.
\]
Substituting these coordinates shows that two such lines are adjacent
exactly in the ten triangles below.  The triangles are edge-disjoint and
there are \(30\) displayed edges.  Since each of the \(15\) lines appears
in four displayed edges, this also proves that
\(\Gamma[\mathcal N_{0}(L_{0})]\) is \(4\)-regular and that no edge is missing:
\begin{equation}\label{eq:H4-N0-triangles}
\begin{array}{lll}
Q_{1}=\{E_{1},X_{+,+},X_{-,-}\},&
Q_{2}=\{E_{1},X_{+,-},X_{-,+}\},\\
Q_{3}=\{E_{2},Y_{+,+},Y_{-,+}\},&
Q_{4}=\{E_{2},Y_{-,-},Y_{+,-}\},\\
Q_{5}=\{E_{3},Z_{+,+},Z_{+,-}\},&
Q_{6}=\{E_{3},Z_{-,+},Z_{-,-}\},\\
Q_{7}=\{X_{+,+},Y_{-,-},Z_{+,-}\},&
Q_{8}=\{Z_{+,+},X_{+,-},Y_{-,+}\},\\
Q_{9}=\{Y_{+,+},Z_{-,+},X_{-,-}\},&
Q_{10}=\{X_{-,+},Z_{-,-},Y_{+,-}\}.
\end{array}
\end{equation}

Each triangle \(Q\) in \eqref{eq:H4-N0-triangles} extends uniquely to an
\(A_{4}\)-subsystem containing \(L_{0}\).  Namely, define
\[
    R(Q)=\{U\in\mathcal L:U\sim L_{0}
    \text{ and }|\{W\in Q:U\sim W\}|=2\},
\]
and set \(C(Q)=\{L_{0}\}\cup Q\cup R(Q)\).  For each of the ten triangles,
the same substitution of coordinates gives
\begin{equation}\label{eq:H4-extension-count}
\begin{array}{c|cccc}
r&0&1&2&3\\
\hline
\#\{U\sim L_{0}:|\{W\in Q:U\sim W\}|=r\}
&2&12&6&0 .
\end{array}
\end{equation}
Thus \(|C(Q)|=10\).  The induced graph \(\Gamma[C(Q)]\) is \(6\)-regular,
two adjacent vertices have three common neighbors, and two non-adjacent
vertices have four common neighbors.  These three assertions are read
directly from the same table of adjacencies used above.  Hence
\(\Gamma[C(Q)]\) is strongly regular with parameters \((10,6,3,4)\), so its
complement is the Petersen graph and \(\Gamma[C(Q)]\cong T(5)\).  Thus
\(C(Q)\) is an \(A_{4}\)-subsystem.

Conversely, if \(C\in\mathcal C\) contains \(L_{0}\), then in the triangular
graph \(T(5)\) the three vertices non-adjacent to a fixed vertex form a
triangle.  Hence \(C\cap\mathcal N_{0}(L_{0})\) must be one of the ten
triangles in \eqref{eq:H4-N0-triangles}, and the remaining six lines are
forced by the definition of \(R(Q)\).  Therefore \(L_{0}\) lies in exactly
\(10\) members of \(\mathcal C\).

The preceding construction and converse are then repeated for the
representatives
\(\left\langle\frac{1}{2}(1,1,1,1)\right\rangle\in\Phi_{2}\) and
\(\langle(0,a,b,c)\rangle\in\Phi_{3}\).  The
orthogonal-neighborhood lists are different, but the same coordinate test
gives exactly ten triangles in \(\Gamma[\mathcal N_{0}(L)]\), and each
triangle extends by the same rule to a unique \(A_{4}\)-subsystem
containing \(L\).  Hence every line lies in exactly \(10\) members of
\(\mathcal C\).  Since every subsystem has \(10\) lines,
double-counting incidences gives \(10|\mathcal C|=10|\mathcal L|=600\),
so \(|\mathcal C|=60\).

It remains to count pair incidences.  We enumerate all pairs
\((L,M)\in\mathcal L^{2}\).  For each fixed line \(L\), list the ten
subsystems through \(L\) by the construction above, and count how often
each second line \(M\) occurs.  The exact count is independent of \(L\) and
depends on \(M\) only through \(|\langle L,M\rangle|\); the multiplicities
are as follows:
\[
\begin{array}{c|ccccc}
|\langle L,M\rangle|
& 1 & \frac{1}{2} & 0 & \frac{\varphi^{-1}}{2} & \frac{\varphi}{2}\\
\hline
\#\{C\in\mathcal C:L,M\in C\}
&10&3&2&0&0 .
\end{array}
\]
Here the entry \(10\) corresponds to \(M=L\).  The entries \(3\) and \(2\)
are also consistent with the internal structure of \(T(5)\): in each
subsystem, a fixed vertex has six adjacent vertices and three non-adjacent
vertices, and the coordinate count shows that these incidences are
distributed uniformly among the \(20\) adjacent and \(15\) orthogonal lines
to \(L\).  Lines with absolute inner product \(\frac{\varphi^{-1}}{2}\) or
\(\frac{\varphi}{2}\) cannot occur with \(L\) in an \(A_{4}\)-subsystem, since the
line system of \(\mathcal V_{5}\) has only the relations \(1\),
\(\frac{1}{2}\), and \(0\).
This proves \eqref{eq:app-H4-A4-pair-table}.
\end{proof}

\end{document}
